% =============================================================
%  introduction.tex -- body text of the dissertation Introduction
%
%  Draft 2026-10-27. Target: <= 3 pages in caltech_thesis.cls.
%  The heading is NOT in this file: dissertation.tex already supplies
%  \extrachapter{Introduction} and \markboth, so replace the \todo{...}
%  line under them with
%      % =============================================================
%  introduction.tex -- body text of the dissertation Introduction
%
%  Draft 2026-10-27. Target: <= 3 pages in caltech_thesis.cls.
%  The heading is NOT in this file: dissertation.tex already supplies
%  \extrachapter{Introduction} and \markboth, so replace the \todo{...}
%  line under them with
%      % =============================================================
%  introduction.tex -- body text of the dissertation Introduction
%
%  Draft 2026-10-27. Target: <= 3 pages in caltech_thesis.cls.
%  The heading is NOT in this file: dissertation.tex already supplies
%  \extrachapter{Introduction} and \markboth, so replace the \todo{...}
%  line under them with
%      % =============================================================
%  introduction.tex -- body text of the dissertation Introduction
%
%  Draft 2026-10-27. Target: <= 3 pages in caltech_thesis.cls.
%  The heading is NOT in this file: dissertation.tex already supplies
%  \extrachapter{Introduction} and \markboth, so replace the \todo{...}
%  line under them with
%      \input{Introduction/introduction}
%
%  NOTE: Do NOT include \documentclass or \begin{document} here.
%
%  New bib keys used (added to references.bib 2026-10-27):
%    iea2026, sovacool2020, esdale2003, dmitrijeva2019
%  Everything else reuses keys already in references.bib.
%
%  Sources for the numbers quoted below:
%    iea2026   -- Global Critical Minerals Outlook 2026 (rev. July 2026):
%                 STEPS demand growth 2024-2040 (p. 114), copper supply
%                 gap ~25% in 2035 (pp. 7, 134), exploration spending
%                 down >10% in 2025 (p. 8), servers/data storage
%                 materials (p. 203), by-product minor minerals (p. 9),
%                 South Africa ~70% of mined platinum in 2025 (p. 193).
%    sovacool2020 -- "geologically determined, yet socially mediated"
%                 (p. 33); PGM substitution difficulty (p. 32).
%    esdale2003 / dmitrijeva2019 -- see the Olympic Dam paragraph.
% =============================================================

Modern energy and computing infrastructure is built from metals whose supply is
concentrated in a select few places and depends on deposits that are difficult to find.
The International Energy Agency projects that, under stated policies, global
demand between 2024 and 2040 will grow by more than a factor of three for
lithium, a factor of two for graphite, 65\% for nickel, about 50\% for magnet rare
earth elements, and more than 25\% for copper, which alone adds roughly 7~million
tons \citep{iea2026}. Two critical innovations are driving this demand: computational infrastructure
in ``the age of AI'', as well as the transition to renewable energy sources, both of which depend on copper, silver, gold and
platinum-group metals, as well as other rare metals such as gallium, germanium and
indium. Technological solutions to the anthropogenic climate crisis
(electric vehicles, battery storage, solar and wind generation, power-grid modernization) 
in effect shift much of the environmental burden from fossil fuel consumption to other, similarly non-renewable
mined materials. Gracefully navigating this transition will require interdisciplinary innovations within the fields of 
geology, chemistry, and biology. Improved techniques in metal extraction from ores through bioengineered fermentation reactors, for example, 
could reduce the environmental footprint of mining and processing while increasing the supply of critical metals \citep{fernando2024}.

Policy plays an essential role in ensuring that the transition to a low-carbon economy is sufficiently rapid and equitable \citep{sovacool2020}. 
When I was conducting geological field work in South Africa in 2023, I was struck by how local communities
viewed large-scale mining efforts. In Draaikraal, Mpumalanga, where economic-grade Platinum Group Element (PGE) ores lie underfoot at shallow 
depths \citep{vandermerwe2007}, 
locals overwhelmingly felt positive about the expansion of mining into their village. To them, the mines represented a source of employment 
and economic growth for their community, wherein livelihoods are currently made from small-scale livestock operations. 
But they also expressed disappointment in how few of their neighbors were hired when the nearest PGE mine (Booysendal, 30~km away) had opened. 
In much of Africa, informal small-scale mining (framed either as "illegal" mining or "artisanal" mining, depending on context) can generate revenue 
for local communities who often benefit little from large-scale mining operations \citep{sovacool2020}, although it also is known to fund 
crime syndicates and spur illegal immigration \citep{martin2019}. These small-scale operations are often unregulated and unsafe: a few months
before this was written, more than a dozen of these independent, informal miners were killed while mining PGEs in an abandoned mine in the 
western limb of the Bushveld Complex in a single incident \citep{ngcobo2026}. 
Such accidents are common. Nonetheless, dispersed and decentralized artisanal mining
operations contribute to the global supply of critical metals from sites that cannot be utilized (and are often abandoned by) larger operations, 
thereby producing needed resources with minimal additional environmental impacts and playing an important role in meeting global demand 
\citep{sovacool2020}. Strong, thoughtful, humanist policies may resolve such challenges. 

I often argue that policy is the most important field of work to direct humanity to a sustainable future, but I also find policy 
intolerably uninteresting and am often discouraged by sluggish political action and the recent regressions within environmental 
policy in the United States and elsewhere. I much prefer to study the natural world. Indeed, my interest in critical metals is rooted 
in the desire to better understand the natural processes that operate within the planet's interior. The application of these studies to 
mineral prospecting to fuel the transition to a low-carbon economy is serendipitous. 

The role of geoscientists to this end is to locate accessible, high-grade deposits of critical metals. Intuitively, mining higher grade deposits 
damages less land and produces less waste to produce a given amount of critical metal. The foundation of this work is to understand
the processes by which these deposits form, change, and survive.

In this thesis I develop geochemical tools for three families of processes wherein natural fluids 
concentrate and rework rare and valuable metals: magmas, metasomatism, and
meteoric mobilization. Magmas deliver critical metals from the mantle to the crust, where they can
become concentrated into cumulates and immiscible sulfide or oxide liquids as they crystallize. Metasomatic 
and hydrothermal fluids then transport and redeposit them, commonly as sulfide minerals. Meteoric (surface) water weathers the
resulting deposits, upgrading some into supergene and lateritic ores and dispersing
others. Each stage leaves a record in minerals, which can be read and interpreted through
geochronometry, which bears on the rates and orders of these processes, and thermodynamic models, which disentangle the chemical behaviors 
of these systems and the effects of the processes that operate on them. Together, these tools allow us to follow the journey of critical metals
within magma through the mantle and hydrothermal reworking into the crust where they can be observed and collected.

The Olympic Dam Iron Oxide Copper Gold (IOCG) deposit in South Australia offers a striking example of why understanding process matters
for discovery. The deposit lies 520~km north-northwest of Adelaide beneath about
260~m of flat-lying sedimentary rock and has no surface expression. It was found
in 1975 by drilling coincident gravity and magnetic anomalies that had been
selected using a conceptual model of copper source, transport and deposition,
together with continental tectonic and lineament analysis. It
holds a resource of 2660~Mt at 1.2\% Cu, 1.4~kg/t U\textsubscript{3}O\textsubscript{8}
and 0.5~g/t Au, which made it the world's largest uranium resource and fifth
largest copper resource at the end of 2001 \citep{esdale2003}. As scientists have 
learned more about the Olympic Dam, more IOCG deposits have been discovered, and the group's growing
diversity challenges the unifying models that explain the origin of IOCG deposits. Nature is complex, and
demands geochemical tools that can be confidently applied and interpreted to disentangle the
processes that produce these resources.

U-Th/He dating is a powerful tool for understanding the timing and duration of geological processes, and
its application to different minerals is sensitive to processes occurring at a variety of temperatures and depths,
including shallow crustal processes that are difficult to study with other geochronometers \citep{farley2002}, making it uniquely
well-suited for studying the full history of an ore deposit, including hydrothermal processes \citep{mcinnes2005}. 
The key to applying U-Th/He dating to a given mineral is to understand how helium behaves within it; this is
the basis on which scientists relate these observations to natural phenomena. For apatite, titanite, and zircon, helium's
behavior at different temperatures is well understood \citep{farley2000,reinersfarley1999,guenthner2013}. More recently, helium in hematite 
has become well understood, expanding the domain of processes accessible to U-Th/He dating \citep{farleyflowers2012,farley2018}. Another 
important target mineral for U-Th/He dating is goethite, a hydrous iron oxide that is a direct chronometer of weathering and laterite formation
\citep{shustervasconcelos2005,vasconcelos2013,monteiro2014}. However, these data are difficult to interpret because goethite decomposes during
conventional helium step-heating experiments, producing data that reflect the kinetics of its decomposition and not helium's behavior within
it in nature \citep{farley2024}. Sulfide minerals, too, are promising targets for U-Th/He dating, because they may directly date primary hydrothermal
mineralization itself \citep{yakubovich2020}, but they too decompose during conventional helium step-heating experiments. 
In Chapter~\ref{ch:apparatus}, I address this problem by developing a method to stabilize these minerals during helium step heat
experiments by sealing samples in silica glass ampoules that maintain water- or sulfur-saturated conditions internally while remaining
permeable to helium. Therein I develop the mathematical and experimental framework to separate the process of helium permeation through the glass
from release by the sample. Chapter~\ref{ch:goethite} applies this method to goethite to measure helium diffusion, demonstrating that it has remarkably high helium retentivity and relatively simple helium systematics. 
Chapter~\ref{ch:sulfide} applies the method to sulfides, the primary hosts of many critical metals. Pyrite and arsenopyrite
prove highly retentive, so much so that pyrite is amongst the most helium-retentive minerals studied, comparable to or exceeding monazite \citep{farley2007}. I then apply the U-Th/He method to several 
sulfide samples to date Mesozoic mineralization at the Gympie gold field in Queensland, Australia, and Archean mineralization
at the Kidd Creek volcanogenic massive sulfide deposit in Ontario, Canada. These results are promising, but expose the complication that 
the uranium and thorium in sulfides are held by inclusions of other minerals, requiring that the interpretation of such data be bespoke
by locality.

Chapter~\ref{ch:ngibbs} turns from clocks to models. Thermodynamic models of
equilibrium phase assemblages, such as \textsc{MELTS} \citep{ghiorsosack1995,gualda2012} and \textsc{HeFESTo} \citep{stixrude2011,stixrude2024},
are the standard way to follow magmas and the mantle through pressure, temperature and composition, but they
are too slow and unstable to be applied to detailed models that more accurately represent the complexity of natural systems 
in which composition varies in space and time.
\emph{nGibbs} is a set of general deep-learning emulators of these models that run
orders of magnitude faster, conserve mass, and can stand in for the parent models
in existing workflows \citep{gleeson2025}. 

Chapter~\ref{ch:draaikraal} concludes the thesis with a field study of what I would consider to be the most important ore deposit in the world: 
the Bushveld Complex in South Africa. Not only does the Bushveld provide the majority of the world's platinum-group metals, as well as 
vast reserves of chromium, vanadium, iron, and titanium, but it is also the type locality for layered mafic intrusions: a 
remarkably well-preserved natural laboratory for studying the fundamental processes that operate within the magmatic plumbing 
systems that supply volcanoes on the surface. At the Bushveld Complex, I studied the platiniferous intersection of the 
Main and Critical Zones with the sedimentary country rock that the Bushveld intrudes, transforms, and consumes. With detailed field work and structural
mapping in the vicinity of the village of Draaikraal in the Mpumalanga Province, I show that the local deformation of the country rock was 
synmagmatic, whereas previous researchers had suggested either that this deformation occurred beforehand \citep{vandermerwe2007}, or 
afterwards as the intrusion cooled and sank \citep{sharpe1982}. I show that this deformation is structurally related to a prominent topographic lineament in the area, and I correlate distinctive marker 
beds across it to show that this feature did not significantly displace rocks on either side of it, arguing instead that this feature coincides
with the edge of the underlying mid-lower Critical Zone, which had buffered the overlying Main Zone to the north from crustal contamination, 
leaving the south to be contaminated by iron-rich fluids that facilitated its oxidation and weathering. 

The Bushveld is a fitting place to end. Indeed, the valley in which Draaikraal resides is blanketed by goethite laterites, lying 
just above the mineralized upper Critical Zone which hosts its platinum group elements within exotic sulfides. Furthermore, the Bushveld 
Complex hosts countless classical problems in igneous petrology that a reliable, fast surrogate thermodynamic model may be applied to, 
exploring the chemical consequences of many scenarios of magma evolution and emplacement to explain the origin of the Bushveld
and other PGE-bearing ore deposits.
%
%  NOTE: Do NOT include \documentclass or \begin{document} here.
%
%  New bib keys used (added to references.bib 2026-10-27):
%    iea2026, sovacool2020, esdale2003, dmitrijeva2019
%  Everything else reuses keys already in references.bib.
%
%  Sources for the numbers quoted below:
%    iea2026   -- Global Critical Minerals Outlook 2026 (rev. July 2026):
%                 STEPS demand growth 2024-2040 (p. 114), copper supply
%                 gap ~25% in 2035 (pp. 7, 134), exploration spending
%                 down >10% in 2025 (p. 8), servers/data storage
%                 materials (p. 203), by-product minor minerals (p. 9),
%                 South Africa ~70% of mined platinum in 2025 (p. 193).
%    sovacool2020 -- "geologically determined, yet socially mediated"
%                 (p. 33); PGM substitution difficulty (p. 32).
%    esdale2003 / dmitrijeva2019 -- see the Olympic Dam paragraph.
% =============================================================

Modern energy and computing infrastructure is built from metals whose supply is
concentrated in a select few places and depends on deposits that are difficult to find.
The International Energy Agency projects that, under stated policies, global
demand between 2024 and 2040 will grow by more than a factor of three for
lithium, a factor of two for graphite, 65\% for nickel, about 50\% for magnet rare
earth elements, and more than 25\% for copper, which alone adds roughly 7~million
tons \citep{iea2026}. Two critical innovations are driving this demand: computational infrastructure
in ``the age of AI'', as well as the transition to renewable energy sources, both of which depend on copper, silver, gold and
platinum-group metals, as well as other rare metals such as gallium, germanium and
indium. Technological solutions to the anthropogenic climate crisis
(electric vehicles, battery storage, solar and wind generation, power-grid modernization) 
in effect shift much of the environmental burden from fossil fuel consumption to other, similarly non-renewable
mined materials. Gracefully navigating this transition will require interdisciplinary innovations within the fields of 
geology, chemistry, and biology. Improved techniques in metal extraction from ores through bioengineered fermentation reactors, for example, 
could reduce the environmental footprint of mining and processing while increasing the supply of critical metals \citep{fernando2024}.

Policy plays an essential role in ensuring that the transition to a low-carbon economy is sufficiently rapid and equitable \citep{sovacool2020}. 
When I was conducting geological field work in South Africa in 2023, I was struck by how local communities
viewed large-scale mining efforts. In Draaikraal, Mpumalanga, where economic-grade Platinum Group Element (PGE) ores lie underfoot at shallow 
depths \citep{vandermerwe2007}, 
locals overwhelmingly felt positive about the expansion of mining into their village. To them, the mines represented a source of employment 
and economic growth for their community, wherein livelihoods are currently made from small-scale livestock operations. 
But they also expressed disappointment in how few of their neighbors were hired when the nearest PGE mine (Booysendal, 30~km away) had opened. 
In much of Africa, informal small-scale mining (framed either as "illegal" mining or "artisanal" mining, depending on context) can generate revenue 
for local communities who often benefit little from large-scale mining operations \citep{sovacool2020}, although it also is known to fund 
crime syndicates and spur illegal immigration \citep{martin2019}. These small-scale operations are often unregulated and unsafe: a few months
before this was written, more than a dozen of these independent, informal miners were killed while mining PGEs in an abandoned mine in the 
western limb of the Bushveld Complex in a single incident \citep{ngcobo2026}. 
Such accidents are common. Nonetheless, dispersed and decentralized artisanal mining
operations contribute to the global supply of critical metals from sites that cannot be utilized (and are often abandoned by) larger operations, 
thereby producing needed resources with minimal additional environmental impacts and playing an important role in meeting global demand 
\citep{sovacool2020}. Strong, thoughtful, humanist policies may resolve such challenges. 

I often argue that policy is the most important field of work to direct humanity to a sustainable future, but I also find policy 
intolerably uninteresting and am often discouraged by sluggish political action and the recent regressions within environmental 
policy in the United States and elsewhere. I much prefer to study the natural world. Indeed, my interest in critical metals is rooted 
in the desire to better understand the natural processes that operate within the planet's interior. The application of these studies to 
mineral prospecting to fuel the transition to a low-carbon economy is serendipitous. 

The role of geoscientists to this end is to locate accessible, high-grade deposits of critical metals. Intuitively, mining higher grade deposits 
damages less land and produces less waste to produce a given amount of critical metal. The foundation of this work is to understand
the processes by which these deposits form, change, and survive.

In this thesis I develop geochemical tools for three families of processes wherein natural fluids 
concentrate and rework rare and valuable metals: magmas, metasomatism, and
meteoric mobilization. Magmas deliver critical metals from the mantle to the crust, where they can
become concentrated into cumulates and immiscible sulfide or oxide liquids as they crystallize. Metasomatic 
and hydrothermal fluids then transport and redeposit them, commonly as sulfide minerals. Meteoric (surface) water weathers the
resulting deposits, upgrading some into supergene and lateritic ores and dispersing
others. Each stage leaves a record in minerals, which can be read and interpreted through
geochronometry, which bears on the rates and orders of these processes, and thermodynamic models, which disentangle the chemical behaviors 
of these systems and the effects of the processes that operate on them. Together, these tools allow us to follow the journey of critical metals
within magma through the mantle and hydrothermal reworking into the crust where they can be observed and collected.

The Olympic Dam Iron Oxide Copper Gold (IOCG) deposit in South Australia offers a striking example of why understanding process matters
for discovery. The deposit lies 520~km north-northwest of Adelaide beneath about
260~m of flat-lying sedimentary rock and has no surface expression. It was found
in 1975 by drilling coincident gravity and magnetic anomalies that had been
selected using a conceptual model of copper source, transport and deposition,
together with continental tectonic and lineament analysis. It
holds a resource of 2660~Mt at 1.2\% Cu, 1.4~kg/t U\textsubscript{3}O\textsubscript{8}
and 0.5~g/t Au, which made it the world's largest uranium resource and fifth
largest copper resource at the end of 2001 \citep{esdale2003}. As scientists have 
learned more about the Olympic Dam, more IOCG deposits have been discovered, and the group's growing
diversity challenges the unifying models that explain the origin of IOCG deposits. Nature is complex, and
demands geochemical tools that can be confidently applied and interpreted to disentangle the
processes that produce these resources.

U-Th/He dating is a powerful tool for understanding the timing and duration of geological processes, and
its application to different minerals is sensitive to processes occurring at a variety of temperatures and depths,
including shallow crustal processes that are difficult to study with other geochronometers \citep{farley2002}, making it uniquely
well-suited for studying the full history of an ore deposit, including hydrothermal processes \citep{mcinnes2005}. 
The key to applying U-Th/He dating to a given mineral is to understand how helium behaves within it; this is
the basis on which scientists relate these observations to natural phenomena. For apatite, titanite, and zircon, helium's
behavior at different temperatures is well understood \citep{farley2000,reinersfarley1999,guenthner2013}. More recently, helium in hematite 
has become well understood, expanding the domain of processes accessible to U-Th/He dating \citep{farleyflowers2012,farley2018}. Another 
important target mineral for U-Th/He dating is goethite, a hydrous iron oxide that is a direct chronometer of weathering and laterite formation
\citep{shustervasconcelos2005,vasconcelos2013,monteiro2014}. However, these data are difficult to interpret because goethite decomposes during
conventional helium step-heating experiments, producing data that reflect the kinetics of its decomposition and not helium's behavior within
it in nature \citep{farley2024}. Sulfide minerals, too, are promising targets for U-Th/He dating, because they may directly date primary hydrothermal
mineralization itself \citep{yakubovich2020}, but they too decompose during conventional helium step-heating experiments. 
In Chapter~\ref{ch:apparatus}, I address this problem by developing a method to stabilize these minerals during helium step heat
experiments by sealing samples in silica glass ampoules that maintain water- or sulfur-saturated conditions internally while remaining
permeable to helium. Therein I develop the mathematical and experimental framework to separate the process of helium permeation through the glass
from release by the sample. Chapter~\ref{ch:goethite} applies this method to goethite to measure helium diffusion, demonstrating that it has remarkably high helium retentivity and relatively simple helium systematics. 
Chapter~\ref{ch:sulfide} applies the method to sulfides, the primary hosts of many critical metals. Pyrite and arsenopyrite
prove highly retentive, so much so that pyrite is amongst the most helium-retentive minerals studied, comparable to or exceeding monazite \citep{farley2007}. I then apply the U-Th/He method to several 
sulfide samples to date Mesozoic mineralization at the Gympie gold field in Queensland, Australia, and Archean mineralization
at the Kidd Creek volcanogenic massive sulfide deposit in Ontario, Canada. These results are promising, but expose the complication that 
the uranium and thorium in sulfides are held by inclusions of other minerals, requiring that the interpretation of such data be bespoke
by locality.

Chapter~\ref{ch:ngibbs} turns from clocks to models. Thermodynamic models of
equilibrium phase assemblages, such as \textsc{MELTS} \citep{ghiorsosack1995,gualda2012} and \textsc{HeFESTo} \citep{stixrude2011,stixrude2024},
are the standard way to follow magmas and the mantle through pressure, temperature and composition, but they
are too slow and unstable to be applied to detailed models that more accurately represent the complexity of natural systems 
in which composition varies in space and time.
\emph{nGibbs} is a set of general deep-learning emulators of these models that run
orders of magnitude faster, conserve mass, and can stand in for the parent models
in existing workflows \citep{gleeson2025}. 

Chapter~\ref{ch:draaikraal} concludes the thesis with a field study of what I would consider to be the most important ore deposit in the world: 
the Bushveld Complex in South Africa. Not only does the Bushveld provide the majority of the world's platinum-group metals, as well as 
vast reserves of chromium, vanadium, iron, and titanium, but it is also the type locality for layered mafic intrusions: a 
remarkably well-preserved natural laboratory for studying the fundamental processes that operate within the magmatic plumbing 
systems that supply volcanoes on the surface. At the Bushveld Complex, I studied the platiniferous intersection of the 
Main and Critical Zones with the sedimentary country rock that the Bushveld intrudes, transforms, and consumes. With detailed field work and structural
mapping in the vicinity of the village of Draaikraal in the Mpumalanga Province, I show that the local deformation of the country rock was 
synmagmatic, whereas previous researchers had suggested either that this deformation occurred beforehand \citep{vandermerwe2007}, or 
afterwards as the intrusion cooled and sank \citep{sharpe1982}. I show that this deformation is structurally related to a prominent topographic lineament in the area, and I correlate distinctive marker 
beds across it to show that this feature did not significantly displace rocks on either side of it, arguing instead that this feature coincides
with the edge of the underlying mid-lower Critical Zone, which had buffered the overlying Main Zone to the north from crustal contamination, 
leaving the south to be contaminated by iron-rich fluids that facilitated its oxidation and weathering. 

The Bushveld is a fitting place to end. Indeed, the valley in which Draaikraal resides is blanketed by goethite laterites, lying 
just above the mineralized upper Critical Zone which hosts its platinum group elements within exotic sulfides. Furthermore, the Bushveld 
Complex hosts countless classical problems in igneous petrology that a reliable, fast surrogate thermodynamic model may be applied to, 
exploring the chemical consequences of many scenarios of magma evolution and emplacement to explain the origin of the Bushveld
and other PGE-bearing ore deposits.
%
%  NOTE: Do NOT include \documentclass or \begin{document} here.
%
%  New bib keys used (added to references.bib 2026-10-27):
%    iea2026, sovacool2020, esdale2003, dmitrijeva2019
%  Everything else reuses keys already in references.bib.
%
%  Sources for the numbers quoted below:
%    iea2026   -- Global Critical Minerals Outlook 2026 (rev. July 2026):
%                 STEPS demand growth 2024-2040 (p. 114), copper supply
%                 gap ~25% in 2035 (pp. 7, 134), exploration spending
%                 down >10% in 2025 (p. 8), servers/data storage
%                 materials (p. 203), by-product minor minerals (p. 9),
%                 South Africa ~70% of mined platinum in 2025 (p. 193).
%    sovacool2020 -- "geologically determined, yet socially mediated"
%                 (p. 33); PGM substitution difficulty (p. 32).
%    esdale2003 / dmitrijeva2019 -- see the Olympic Dam paragraph.
% =============================================================

Modern energy and computing infrastructure is built from metals whose supply is
concentrated in a select few places and depends on deposits that are difficult to find.
The International Energy Agency projects that, under stated policies, global
demand between 2024 and 2040 will grow by more than a factor of three for
lithium, a factor of two for graphite, 65\% for nickel, about 50\% for magnet rare
earth elements, and more than 25\% for copper, which alone adds roughly 7~million
tons \citep{iea2026}. Two critical innovations are driving this demand: computational infrastructure
in ``the age of AI'', as well as the transition to renewable energy sources, both of which depend on copper, silver, gold and
platinum-group metals, as well as other rare metals such as gallium, germanium and
indium. Technological solutions to the anthropogenic climate crisis
(electric vehicles, battery storage, solar and wind generation, power-grid modernization) 
in effect shift much of the environmental burden from fossil fuel consumption to other, similarly non-renewable
mined materials. Gracefully navigating this transition will require interdisciplinary innovations within the fields of 
geology, chemistry, and biology. Improved techniques in metal extraction from ores through bioengineered fermentation reactors, for example, 
could reduce the environmental footprint of mining and processing while increasing the supply of critical metals \citep{fernando2024}.

Policy plays an essential role in ensuring that the transition to a low-carbon economy is sufficiently rapid and equitable \citep{sovacool2020}. 
When I was conducting geological field work in South Africa in 2023, I was struck by how local communities
viewed large-scale mining efforts. In Draaikraal, Mpumalanga, where economic-grade Platinum Group Element (PGE) ores lie underfoot at shallow 
depths \citep{vandermerwe2007}, 
locals overwhelmingly felt positive about the expansion of mining into their village. To them, the mines represented a source of employment 
and economic growth for their community, wherein livelihoods are currently made from small-scale livestock operations. 
But they also expressed disappointment in how few of their neighbors were hired when the nearest PGE mine (Booysendal, 30~km away) had opened. 
In much of Africa, informal small-scale mining (framed either as "illegal" mining or "artisanal" mining, depending on context) can generate revenue 
for local communities who often benefit little from large-scale mining operations \citep{sovacool2020}, although it also is known to fund 
crime syndicates and spur illegal immigration \citep{martin2019}. These small-scale operations are often unregulated and unsafe: a few months
before this was written, more than a dozen of these independent, informal miners were killed while mining PGEs in an abandoned mine in the 
western limb of the Bushveld Complex in a single incident \citep{ngcobo2026}. 
Such accidents are common. Nonetheless, dispersed and decentralized artisanal mining
operations contribute to the global supply of critical metals from sites that cannot be utilized (and are often abandoned by) larger operations, 
thereby producing needed resources with minimal additional environmental impacts and playing an important role in meeting global demand 
\citep{sovacool2020}. Strong, thoughtful, humanist policies may resolve such challenges. 

I often argue that policy is the most important field of work to direct humanity to a sustainable future, but I also find policy 
intolerably uninteresting and am often discouraged by sluggish political action and the recent regressions within environmental 
policy in the United States and elsewhere. I much prefer to study the natural world. Indeed, my interest in critical metals is rooted 
in the desire to better understand the natural processes that operate within the planet's interior. The application of these studies to 
mineral prospecting to fuel the transition to a low-carbon economy is serendipitous. 

The role of geoscientists to this end is to locate accessible, high-grade deposits of critical metals. Intuitively, mining higher grade deposits 
damages less land and produces less waste to produce a given amount of critical metal. The foundation of this work is to understand
the processes by which these deposits form, change, and survive.

In this thesis I develop geochemical tools for three families of processes wherein natural fluids 
concentrate and rework rare and valuable metals: magmas, metasomatism, and
meteoric mobilization. Magmas deliver critical metals from the mantle to the crust, where they can
become concentrated into cumulates and immiscible sulfide or oxide liquids as they crystallize. Metasomatic 
and hydrothermal fluids then transport and redeposit them, commonly as sulfide minerals. Meteoric (surface) water weathers the
resulting deposits, upgrading some into supergene and lateritic ores and dispersing
others. Each stage leaves a record in minerals, which can be read and interpreted through
geochronometry, which bears on the rates and orders of these processes, and thermodynamic models, which disentangle the chemical behaviors 
of these systems and the effects of the processes that operate on them. Together, these tools allow us to follow the journey of critical metals
within magma through the mantle and hydrothermal reworking into the crust where they can be observed and collected.

The Olympic Dam Iron Oxide Copper Gold (IOCG) deposit in South Australia offers a striking example of why understanding process matters
for discovery. The deposit lies 520~km north-northwest of Adelaide beneath about
260~m of flat-lying sedimentary rock and has no surface expression. It was found
in 1975 by drilling coincident gravity and magnetic anomalies that had been
selected using a conceptual model of copper source, transport and deposition,
together with continental tectonic and lineament analysis. It
holds a resource of 2660~Mt at 1.2\% Cu, 1.4~kg/t U\textsubscript{3}O\textsubscript{8}
and 0.5~g/t Au, which made it the world's largest uranium resource and fifth
largest copper resource at the end of 2001 \citep{esdale2003}. As scientists have 
learned more about the Olympic Dam, more IOCG deposits have been discovered, and the group's growing
diversity challenges the unifying models that explain the origin of IOCG deposits. Nature is complex, and
demands geochemical tools that can be confidently applied and interpreted to disentangle the
processes that produce these resources.

U-Th/He dating is a powerful tool for understanding the timing and duration of geological processes, and
its application to different minerals is sensitive to processes occurring at a variety of temperatures and depths,
including shallow crustal processes that are difficult to study with other geochronometers \citep{farley2002}, making it uniquely
well-suited for studying the full history of an ore deposit, including hydrothermal processes \citep{mcinnes2005}. 
The key to applying U-Th/He dating to a given mineral is to understand how helium behaves within it; this is
the basis on which scientists relate these observations to natural phenomena. For apatite, titanite, and zircon, helium's
behavior at different temperatures is well understood \citep{farley2000,reinersfarley1999,guenthner2013}. More recently, helium in hematite 
has become well understood, expanding the domain of processes accessible to U-Th/He dating \citep{farleyflowers2012,farley2018}. Another 
important target mineral for U-Th/He dating is goethite, a hydrous iron oxide that is a direct chronometer of weathering and laterite formation
\citep{shustervasconcelos2005,vasconcelos2013,monteiro2014}. However, these data are difficult to interpret because goethite decomposes during
conventional helium step-heating experiments, producing data that reflect the kinetics of its decomposition and not helium's behavior within
it in nature \citep{farley2024}. Sulfide minerals, too, are promising targets for U-Th/He dating, because they may directly date primary hydrothermal
mineralization itself \citep{yakubovich2020}, but they too decompose during conventional helium step-heating experiments. 
In Chapter~\ref{ch:apparatus}, I address this problem by developing a method to stabilize these minerals during helium step heat
experiments by sealing samples in silica glass ampoules that maintain water- or sulfur-saturated conditions internally while remaining
permeable to helium. Therein I develop the mathematical and experimental framework to separate the process of helium permeation through the glass
from release by the sample. Chapter~\ref{ch:goethite} applies this method to goethite to measure helium diffusion, demonstrating that it has remarkably high helium retentivity and relatively simple helium systematics. 
Chapter~\ref{ch:sulfide} applies the method to sulfides, the primary hosts of many critical metals. Pyrite and arsenopyrite
prove highly retentive, so much so that pyrite is amongst the most helium-retentive minerals studied, comparable to or exceeding monazite \citep{farley2007}. I then apply the U-Th/He method to several 
sulfide samples to date Mesozoic mineralization at the Gympie gold field in Queensland, Australia, and Archean mineralization
at the Kidd Creek volcanogenic massive sulfide deposit in Ontario, Canada. These results are promising, but expose the complication that 
the uranium and thorium in sulfides are held by inclusions of other minerals, requiring that the interpretation of such data be bespoke
by locality.

Chapter~\ref{ch:ngibbs} turns from clocks to models. Thermodynamic models of
equilibrium phase assemblages, such as \textsc{MELTS} \citep{ghiorsosack1995,gualda2012} and \textsc{HeFESTo} \citep{stixrude2011,stixrude2024},
are the standard way to follow magmas and the mantle through pressure, temperature and composition, but they
are too slow and unstable to be applied to detailed models that more accurately represent the complexity of natural systems 
in which composition varies in space and time.
\emph{nGibbs} is a set of general deep-learning emulators of these models that run
orders of magnitude faster, conserve mass, and can stand in for the parent models
in existing workflows \citep{gleeson2025}. 

Chapter~\ref{ch:draaikraal} concludes the thesis with a field study of what I would consider to be the most important ore deposit in the world: 
the Bushveld Complex in South Africa. Not only does the Bushveld provide the majority of the world's platinum-group metals, as well as 
vast reserves of chromium, vanadium, iron, and titanium, but it is also the type locality for layered mafic intrusions: a 
remarkably well-preserved natural laboratory for studying the fundamental processes that operate within the magmatic plumbing 
systems that supply volcanoes on the surface. At the Bushveld Complex, I studied the platiniferous intersection of the 
Main and Critical Zones with the sedimentary country rock that the Bushveld intrudes, transforms, and consumes. With detailed field work and structural
mapping in the vicinity of the village of Draaikraal in the Mpumalanga Province, I show that the local deformation of the country rock was 
synmagmatic, whereas previous researchers had suggested either that this deformation occurred beforehand \citep{vandermerwe2007}, or 
afterwards as the intrusion cooled and sank \citep{sharpe1982}. I show that this deformation is structurally related to a prominent topographic lineament in the area, and I correlate distinctive marker 
beds across it to show that this feature did not significantly displace rocks on either side of it, arguing instead that this feature coincides
with the edge of the underlying mid-lower Critical Zone, which had buffered the overlying Main Zone to the north from crustal contamination, 
leaving the south to be contaminated by iron-rich fluids that facilitated its oxidation and weathering. 

The Bushveld is a fitting place to end. Indeed, the valley in which Draaikraal resides is blanketed by goethite laterites, lying 
just above the mineralized upper Critical Zone which hosts its platinum group elements within exotic sulfides. Furthermore, the Bushveld 
Complex hosts countless classical problems in igneous petrology that a reliable, fast surrogate thermodynamic model may be applied to, 
exploring the chemical consequences of many scenarios of magma evolution and emplacement to explain the origin of the Bushveld
and other PGE-bearing ore deposits.
%
%  NOTE: Do NOT include \documentclass or \begin{document} here.
%
%  New bib keys used (added to references.bib 2026-10-27):
%    iea2026, sovacool2020, esdale2003, dmitrijeva2019
%  Everything else reuses keys already in references.bib.
%
%  Sources for the numbers quoted below:
%    iea2026   -- Global Critical Minerals Outlook 2026 (rev. July 2026):
%                 STEPS demand growth 2024-2040 (p. 114), copper supply
%                 gap ~25% in 2035 (pp. 7, 134), exploration spending
%                 down >10% in 2025 (p. 8), servers/data storage
%                 materials (p. 203), by-product minor minerals (p. 9),
%                 South Africa ~70% of mined platinum in 2025 (p. 193).
%    sovacool2020 -- "geologically determined, yet socially mediated"
%                 (p. 33); PGM substitution difficulty (p. 32).
%    esdale2003 / dmitrijeva2019 -- see the Olympic Dam paragraph.
% =============================================================

Modern energy and computing infrastructure is built from metals whose supply is
concentrated in a select few places and depends on deposits that are difficult to find.
The International Energy Agency projects that, under stated policies, global
demand between 2024 and 2040 will grow by more than a factor of three for
lithium, a factor of two for graphite, 65\% for nickel, about 50\% for magnet rare
earth elements, and more than 25\% for copper, which alone adds roughly 7~million
tons \citep{iea2026}. Two critical innovations are driving this demand: computational infrastructure
in ``the age of AI'', as well as the transition to renewable energy sources, both of which depend on copper, silver, gold and
platinum-group metals, as well as other rare metals such as gallium, germanium and
indium. Technological solutions to the anthropogenic climate crisis
(electric vehicles, battery storage, solar and wind generation, power-grid modernization) 
in effect shift much of the environmental burden from fossil fuel consumption to other, similarly non-renewable
mined materials. Gracefully navigating this transition will require interdisciplinary innovations within the fields of 
geology, chemistry, and biology. Improved techniques in metal extraction from ores through bioengineered fermentation reactors, for example, 
could reduce the environmental footprint of mining and processing while increasing the supply of critical metals \citep{fernando2024}.

Policy plays an essential role in ensuring that the transition to a low-carbon economy is sufficiently rapid and equitable \citep{sovacool2020}. 
When I was conducting geological field work in South Africa in 2023, I was struck by how local communities
viewed large-scale mining efforts. In Draaikraal, Mpumalanga, where economic-grade Platinum Group Element (PGE) ores lie underfoot at shallow 
depths \citep{vandermerwe2007}, 
locals overwhelmingly felt positive about the expansion of mining into their village. To them, the mines represented a source of employment 
and economic growth for their community, wherein livelihoods are currently made from small-scale livestock operations. 
But they also expressed disappointment in how few of their neighbors were hired when the nearest PGE mine (Booysendal, 30~km away) had opened. 
In much of Africa, informal small-scale mining (framed either as "illegal" mining or "artisanal" mining, depending on context) can generate revenue 
for local communities who often benefit little from large-scale mining operations \citep{sovacool2020}, although it also is known to fund 
crime syndicates and spur illegal immigration \citep{martin2019}. These small-scale operations are often unregulated and unsafe: a few months
before this was written, more than a dozen of these independent, informal miners were killed while mining PGEs in an abandoned mine in the 
western limb of the Bushveld Complex in a single incident \citep{ngcobo2026}. 
Such accidents are common. Nonetheless, dispersed and decentralized artisanal mining
operations contribute to the global supply of critical metals from sites that cannot be utilized (and are often abandoned by) larger operations, 
thereby producing needed resources with minimal additional environmental impacts and playing an important role in meeting global demand 
\citep{sovacool2020}. Strong, thoughtful, humanist policies may resolve such challenges. 

I often argue that policy is the most important field of work to direct humanity to a sustainable future, but I also find policy 
intolerably uninteresting and am often discouraged by sluggish political action and the recent regressions within environmental 
policy in the United States and elsewhere. I much prefer to study the natural world. Indeed, my interest in critical metals is rooted 
in the desire to better understand the natural processes that operate within the planet's interior. The application of these studies to 
mineral prospecting to fuel the transition to a low-carbon economy is serendipitous. 

The role of geoscientists to this end is to locate accessible, high-grade deposits of critical metals. Intuitively, mining higher grade deposits 
damages less land and produces less waste to produce a given amount of critical metal. The foundation of this work is to understand
the processes by which these deposits form, change, and survive.

In this thesis I develop geochemical tools for three families of processes wherein natural fluids 
concentrate and rework rare and valuable metals: magmas, metasomatism, and
meteoric mobilization. Magmas deliver critical metals from the mantle to the crust, where they can
become concentrated into cumulates and immiscible sulfide or oxide liquids as they crystallize. Metasomatic 
and hydrothermal fluids then transport and redeposit them, commonly as sulfide minerals. Meteoric (surface) water weathers the
resulting deposits, upgrading some into supergene and lateritic ores and dispersing
others. Each stage leaves a record in minerals, which can be read and interpreted through
geochronometry, which bears on the rates and orders of these processes, and thermodynamic models, which disentangle the chemical behaviors 
of these systems and the effects of the processes that operate on them. Together, these tools allow us to follow the journey of critical metals
within magma through the mantle and hydrothermal reworking into the crust where they can be observed and collected.

The Olympic Dam Iron Oxide Copper Gold (IOCG) deposit in South Australia offers a striking example of why understanding process matters
for discovery. The deposit lies 520~km north-northwest of Adelaide beneath about
260~m of flat-lying sedimentary rock and has no surface expression. It was found
in 1975 by drilling coincident gravity and magnetic anomalies that had been
selected using a conceptual model of copper source, transport and deposition,
together with continental tectonic and lineament analysis. It
holds a resource of 2660~Mt at 1.2\% Cu, 1.4~kg/t U\textsubscript{3}O\textsubscript{8}
and 0.5~g/t Au, which made it the world's largest uranium resource and fifth
largest copper resource at the end of 2001 \citep{esdale2003}. As scientists have 
learned more about the Olympic Dam, more IOCG deposits have been discovered, and the group's growing
diversity challenges the unifying models that explain the origin of IOCG deposits. Nature is complex, and
demands geochemical tools that can be confidently applied and interpreted to disentangle the
processes that produce these resources.

U-Th/He dating is a powerful tool for understanding the timing and duration of geological processes, and
its application to different minerals is sensitive to processes occurring at a variety of temperatures and depths,
including shallow crustal processes that are difficult to study with other geochronometers \citep{farley2002}, making it uniquely
well-suited for studying the full history of an ore deposit, including hydrothermal processes \citep{mcinnes2005}. 
The key to applying U-Th/He dating to a given mineral is to understand how helium behaves within it; this is
the basis on which scientists relate these observations to natural phenomena. For apatite, titanite, and zircon, helium's
behavior at different temperatures is well understood \citep{farley2000,reinersfarley1999,guenthner2013}. More recently, helium in hematite 
has become well understood, expanding the domain of processes accessible to U-Th/He dating \citep{farleyflowers2012,farley2018}. Another 
important target mineral for U-Th/He dating is goethite, a hydrous iron oxide that is a direct chronometer of weathering and laterite formation
\citep{shustervasconcelos2005,vasconcelos2013,monteiro2014}. However, these data are difficult to interpret because goethite decomposes during
conventional helium step-heating experiments, producing data that reflect the kinetics of its decomposition and not helium's behavior within
it in nature \citep{farley2024}. Sulfide minerals, too, are promising targets for U-Th/He dating, because they may directly date primary hydrothermal
mineralization itself \citep{yakubovich2020}, but they too decompose during conventional helium step-heating experiments. 
In Chapter~\ref{ch:apparatus}, I address this problem by developing a method to stabilize these minerals during helium step heat
experiments by sealing samples in silica glass ampoules that maintain water- or sulfur-saturated conditions internally while remaining
permeable to helium. Therein I develop the mathematical and experimental framework to separate the process of helium permeation through the glass
from release by the sample. Chapter~\ref{ch:goethite} applies this method to goethite to measure helium diffusion, demonstrating that it has remarkably high helium retentivity and relatively simple helium systematics. 
Chapter~\ref{ch:sulfide} applies the method to sulfides, the primary hosts of many critical metals. Pyrite and arsenopyrite
prove highly retentive, so much so that pyrite is amongst the most helium-retentive minerals studied, comparable to or exceeding monazite \citep{farley2007}. I then apply the U-Th/He method to several 
sulfide samples to date Mesozoic mineralization at the Gympie gold field in Queensland, Australia, and Archean mineralization
at the Kidd Creek volcanogenic massive sulfide deposit in Ontario, Canada. These results are promising, but expose the complication that 
the uranium and thorium in sulfides are held by inclusions of other minerals, requiring that the interpretation of such data be bespoke
by locality.

Chapter~\ref{ch:ngibbs} turns from clocks to models. Thermodynamic models of
equilibrium phase assemblages, such as \textsc{MELTS} \citep{ghiorsosack1995,gualda2012} and \textsc{HeFESTo} \citep{stixrude2011,stixrude2024},
are the standard way to follow magmas and the mantle through pressure, temperature and composition, but they
are too slow and unstable to be applied to detailed models that more accurately represent the complexity of natural systems 
in which composition varies in space and time.
\emph{nGibbs} is a set of general deep-learning emulators of these models that run
orders of magnitude faster, conserve mass, and can stand in for the parent models
in existing workflows \citep{gleeson2025}. 

Chapter~\ref{ch:draaikraal} concludes the thesis with a field study of what I would consider to be the most important ore deposit in the world: 
the Bushveld Complex in South Africa. Not only does the Bushveld provide the majority of the world's platinum-group metals, as well as 
vast reserves of chromium, vanadium, iron, and titanium, but it is also the type locality for layered mafic intrusions: a 
remarkably well-preserved natural laboratory for studying the fundamental processes that operate within the magmatic plumbing 
systems that supply volcanoes on the surface. At the Bushveld Complex, I studied the platiniferous intersection of the 
Main and Critical Zones with the sedimentary country rock that the Bushveld intrudes, transforms, and consumes. With detailed field work and structural
mapping in the vicinity of the village of Draaikraal in the Mpumalanga Province, I show that the local deformation of the country rock was 
synmagmatic, whereas previous researchers had suggested either that this deformation occurred beforehand \citep{vandermerwe2007}, or 
afterwards as the intrusion cooled and sank \citep{sharpe1982}. I show that this deformation is structurally related to a prominent topographic lineament in the area, and I correlate distinctive marker 
beds across it to show that this feature did not significantly displace rocks on either side of it, arguing instead that this feature coincides
with the edge of the underlying mid-lower Critical Zone, which had buffered the overlying Main Zone to the north from crustal contamination, 
leaving the south to be contaminated by iron-rich fluids that facilitated its oxidation and weathering. 

The Bushveld is a fitting place to end. Indeed, the valley in which Draaikraal resides is blanketed by goethite laterites, lying 
just above the mineralized upper Critical Zone which hosts its platinum group elements within exotic sulfides. Furthermore, the Bushveld 
Complex hosts countless classical problems in igneous petrology that a reliable, fast surrogate thermodynamic model may be applied to, 
exploring the chemical consequences of many scenarios of magma evolution and emplacement to explain the origin of the Bushveld
and other PGE-bearing ore deposits.